\section{Removing the global safety stops}\label{sec:global}

All estimates here are proved for the stopped process. On it,
\(k\le K\), \(r\ge\rho\), and \(\eta\ge1-\epsilon_n\), so
Proposition~\ref{prop:mixed} is available before any probability bound
on the global stop has been established.

\subsection{Pair-correction and virtual-deletion guards}
\begin{lemma}[Pair-correction tail]\label{lem:etatail}
For each fixed positive integer \(m\),
\begin{equation}\label{eq:etatail}
 \Prob(\text{pair-correction guard})
 \le C_m n^{4-m+2m\xi+d(m-1)+m\zeta}.
\end{equation}
\end{lemma}
\begin{proof}
For a fixed pair \(a,b\), classify a cause changing both roots
according to whether the two lines are ordinary or designated.
Two ordinary lines require two distinct selected witnesses; one
ordinary line requires one; two designated lines require only an
intersection of two among finitely many directions.

In an \(m\)-fold expansion, regard distinct selected witnesses as
vertices and the double-ordinary causes as edges of a multigraph.
If there are \(j\) witnesses and \(c\) connected components, the
first edge of a component costs \(C_\xi n^{2+\xi}\) by
\eqref{eq:twostar}. An edge adding one new witness costs
\(C_\xi n^{1+\xi}\) by \eqref{eq:fixedline}. With both witnesses
known, the cause is the intersection of two fixed lines and has at
most one choice. Coincidence of those lines would put the cause on
the excluded baseline. Thus the count is at most
\(C_m n^{j+(j-c)\xi}\le C_m n^{j+2m\xi}\).
A component containing a designated/ordinary cause may be started
with that cause, at cost \(C_\xi n^{1+\xi}\). A double-designated
cause has at most \(g^2\) choices. This also covers mixed components
and a witness appearing on different sides of different edges.

If the physical causes have \(k\ge1\) distinct labels, their rates
supply \(n^{-m}\), root repetitions cost
\(B^{m-k}\le B^{m-1}\), and the \(j\) distinct witnesses cost
\(C_m n^{-j}\) by Proposition~\ref{prop:mixed}. Root activity
indicators can be replaced by constant multiples of root weights
using Lemma~\ref{lem:product}. A label that is both an active cause
and a selected witness gives zero. Therefore
\[
 \E\kappa_{ab}(t)^m\le C_m n^{-m+2m\xi}B^{m-1}.
\]
Hitting the floor requires
\(\int_0^T\kappa_{ab}\dd t\ge-\log(1-\epsilon_n)\ge\epsilon_n\).
Apply Holder to the time integral, then Markov and a union bound
over at most \(n^4\) pairs. This is \eqref{eq:etatail}.
\end{proof}

\begin{lemma}[Virtual-deletion tail]\label{lem:virtualtail}
For each fixed positive integer \(m\),
\begin{equation}\label{eq:virtualtail}
 \Prob(\text{virtual-deletion guard})
 \le C_m\sum_{a=1}^m n^{3-\gamma m+d(m-a)+a\xi}.
\end{equation}
\end{lemma}
\begin{proof}
Fix a designated \(\ell\) and an active \(q\notin\ell\).
In an \(m\)-fold expansion of \eqref{eq:virtual}, suppose there are
\(a\) distinct leaves on \(\ell\). The lines from \(q\) to these
leaves are distinct, and their selected witnesses are distinct:
a common witness on two such lines would be \(q\), which is active.
By \eqref{eq:fixedline} the number of leaf/witness choices is at most
\(C_m n^{a(1+\xi)}\). Source prices supply \(n^{-a}\), and repeated
leaves cost \(B^{m-a}\). The normalized weight is at most \(2w\).
Consequently
\[
 \E[\ind_{q\text{ active}}\mathcal V_{\ell,q}^m]
 \le C_m\sum_{a=1}^m B^{m-a}n^{a\xi}.
\]
Apply the same fixed-label hierarchy at the first crossing; no
union over observation times is necessary. There are \(O(n)\)
designated lines and \(n^2\) possible external points. Markov at
level \(n^\gamma\) gives \eqref{eq:virtualtail}.
\end{proof}

The slope in \(m\) in \eqref{eq:etatail} is
\(-(1-d-2\xi-\zeta)<0\). Since \(\xi<d\), the largest exponent in
\eqref{eq:virtualtail} occurs at \(a=1\) and equals
\(3-d+\xi-m(\gamma-d)\). Both probabilities are therefore smaller
than any prescribed fixed power of \(n^{-1}\), after choosing one
sufficiently large fixed moment order. The process itself is unchanged
by that choice of proof order.

\subsection{The protected law of a designated line}
Fix a designated line and one of its occupancy stages \(N=0,1\).
Starting at the stage's entry time \(U\), suppress its next birth
through the density process
\[
 D_t=\ind_{\text{no new line birth in }(U,t]}
 \exp\left(\int_U^t\Lambda_\ell(s)\dd s\right),
\]
stopped and saved at the relevant guards or \(T\). Before \(U\) its
value is one. The birth at \(U\), if any, starts the stage and is
not counted again as its next birth. The continuous drift
\(D\Lambda_\ell\) cancels the zeroing-jump drift. Also
\[
 \int_U^t\Lambda_\ell\dd s
 =\int_{k(U)}^{k(t)}a_N(k)\dd k
 \le2\log\frac2{2-K}.
\]
Thus \(D\) is a bounded martingale and defines a probability law
conditionally on the entry history. Its upper bound is the fixed
constant \(A_0^2\), where \(A_0=2/(2-K)\). On a crossing occurring
before the next line birth, \(D\ge1\), so the original crossing
probability is at most its probability under the protected law.

Normalize each active line weight by
\(c_p=w_p/a_N(k_\ell)\), and save this \emph{normalized} value at
positive retirement. It is not subsequently divided by a changing
\(a_N\). For a fixed cell/type subset of the line, let \(M\) be the
sum of current and saved \(c_p\). Define its entropy sum using
\[
 h(c)=c\left(1+\log\frac{2B}{c}\right),\quad h(0)=0,
 \qquad H=\sum h(c_p),
\]
again with saved terms frozen. Under the protected law the surviving
positive drift of each credit is
\(g_p=\lambda_p(1-b_N)c_p\). This follows by subtracting
\(\mathrm d\log a_N/\mathrm dt=\Lambda_\ell b_N\) from the original
ODE and deleting only the line-birth jump rates. It is
\(\lambda_pc_p/2\) at occupancy zero and zero at occupancy one.

\begin{lemma}[Credit, entropy and whole-event energy]\label{lem:entropy}
Stop further at \(M\le C_0n\). There are compensated martingales
\(Y_M,Y_H\) and a nonnegative cumulative cost \(V\) such that
\begin{align}
 M&=M(U)+D_M+Y_M,&D_M&=O(B),\label{eq:creditdecomp}\\
 H&=H(U)+D_H-V+Y_H,&D_H&=O(B\log n),\label{eq:entropydecomp}
\end{align}
and, before the virtual guard,
\begin{align}
 \langle Y_M\rangle&\le C n^\gamma V,
 &|\Delta Y_M|&\le C n^\gamma,\label{eq:bracketM}\\
 \langle Y_H\rangle&\le C n^\gamma(\log n)^2V,
 &|\Delta Y_H|&\le C n^\gamma\log n.\label{eq:bracketH}
\end{align}
The constants are independent of the admissible entry history.
\end{lemma}
\begin{proof}
For a multiplier \(J\in\{0,1/2,1\}\), direct algebra gives
\[
 h'(c)c(1-J)+h(Jc)-h(c)=-c(J\log J-J+1).
\]
The cost is \(c\) for deletion, \(c(1-\log2)/2\) for halving,
and zero otherwise. It is comparable by fixed positive constants
to the credit decrease \(c(1-J)\). Define \(V\) to be the rate
integral of the \emph{sum of these costs for each complete birth
event}. Compensating gives \eqref{eq:creditdecomp} and
\eqref{eq:entropydecomp}. Since \(\lambda_p\le B/n\) and the
credit sum is capped by \(C_0n\), the first drift is \(O(B)\).
The active credits have the positive lower bound supplied by
Lemma~\ref{lem:product} divided by \(A_0\), so
\(h'(c)=\log(2B/c)=O(\log n)\) on them. This gives the entropy
drift estimate.

In one external birth event, the ordinary-line deletion of normalized
mass is at most \(M_n\), and effects through other designated
directions touch at most \(g\) points, each of credit at most \(2B\).
Its total credit decrease \(\Delta\) is therefore at most
\(M_n+2gB=O(n^\gamma)\). The square of the \emph{whole} decrease
satisfies \(\Delta^2\le C n^\gamma\Delta\).
The absolute entropy decrease is at most \(C\log n\) times this
credit decrease, including deletion to zero. Comparing the decrease
to the event's total cost and integrating its birth rate proves
\eqref{eq:bracketM}--\eqref{eq:bracketH}. These estimates retain
all cross terms of simultaneous changes. Summing the squares of
individual point changes would not prove the required bracket bound.
\end{proof}

\begin{lemma}[Exponential crossing estimate]\label{lem:exp}
Let \(Y\) be a compensated birth martingale with \(Y_0=0\), jumps
bounded in absolute value by \(b\), and continuous predictable
quadratic variation \(Q\). For \(x,v>0\),
\begin{equation}\label{eq:freedman}
 \Prob(\exists t\le T:Y_t\ge x,\ Q_t\le v)
 \le\exp\left(-\frac{x^2}{2(v+bx/3)}\right).
\end{equation}
\end{lemma}
\begin{proof}
For \(0<\theta b<3\) and \(|u|\le b\), the Taylor series and
\(k!\ge2\cdot3^{k-2}\) for \(k\ge2\) give
\[
 e^{\theta u}-1-\theta u
 \le\psi u^2,\qquad \psi=\frac{\theta^2}{2(1-\theta b/3)}.
\]
Stop at the first \(Y\ge x\), capped at \(T\). The stopped
exponential \(\exp(\theta Y-\psi Q)\) is bounded by
\(\exp(\theta(x+b))\). Its compensated drift is nonpositive by the
displayed inequality, so its expectation is at most one.
Since \(Q\) is continuous and nondecreasing, a crossing with
\(Q\le v\) exists precisely when the first crossing has \(Q\le v\).
On this event the terminal stopped exponential is at least
\(\exp(\theta x-\psi v)\). Markov gives
\(\exp(-\theta x+\psi v)\); take \(\theta=x/(v+bx/3)\).
This proof uses a bounded stopped exponential, rather than assuming
that an unlocalized exponential local martingale has mean one.
\end{proof}

\begin{proposition}[Clock stops are rare]\label{prop:clock}
For each fixed positive tolerance and admissible fixed data, the
failure probability of the simultaneous designated-line credit
bound is at most
\[
 Cn\exp\left(-c\frac{n^{1-\gamma}}{\log n}\right).
\]
In particular, the probability of a designated clock reaching \(K\)
is smaller than every fixed inverse power of \(n\).
\end{proposition}
\begin{proof}
Stop the cumulative cost at \(C_1n\log n\). Since \(H\ge0\),
\(H(U)=O(n\log n)\), and \(D_H=o(n\log n)\), reaching this cost
cap, for sufficiently large \(C_1\), forces an upward fluctuation
of \(Y_H\) of order \(n\log n\). Apply Lemma~\ref{lem:exp} with
\(v=O(n^{1+\gamma}(\log n)^3)\) and
\(b=O(n^\gamma\log n)\). An increase of \(M\) by a fixed multiple
of \(n\) similarly forces a fluctuation of \(Y_M\) of order \(n\),
with \(v=O(n^{1+\gamma}\log n)\) and \(b=O(n^\gamma)\).
Both exponents have order \(n^{1-\gamma}/\log n\).

At the first birth on the line,
\((w/2)/a_1=w/a_0\); the selected point is removed. The saved ledger
is unchanged and the moving normalized credit cannot increase.
At the second birth, all remaining moving credit disappears. To
treat the occupancy-one stage, condition on its actual entry history
and the preceding good credit bound. The entropy bound follows
from that credit bound and \(c\le2B\). The same constants therefore
apply to every such history. Average the conditional estimate; do
not take a union over possible histories or starting times.

There are \(O(n)\) designated lines and a fixed number of cell/type
pieces and phases. Summing their crossing probabilities preserves
the displayed bound. On the resulting good event, integrate the
moving normalized credit with the masks to obtain
\[
 k_\ell\le\text{deterministic line capacity}
             +C_{\mathrm{data}}T\varepsilon+O(n^{-1}).
\]
Choose \(\varepsilon>0\) first so that the second term is smaller
than half of \(\kappa_H\), then choose \(n\) large. The capacity
condition \eqref{eq:heavyinput} precludes hitting \(K\).
Positive retirement need not be rare for this argument: its credit
was retained in the ledger throughout.
\end{proof}

Combining the three guards shows, for every fixed \(A_0'>0\),
\begin{equation}\label{eq:globaltail}
 \Prob(\sigma\le T)\le C_{A_0'}n^{-A_0'}.
\end{equation}
Under any point law this is at most \(B\) times as large, hence is
\(o(1)\) uniformly in the root. This establishes global-stop smallness
before the weight-cap and ordinary-line tag estimates below.



